\begin{theorem}[d one]
\uses{d}
\label{d_one}
\lean{ArithmeticFunction.d_one}
% at /home/runner/work/trough/trough/library/PrimeNumberTheoremAnd/IwaniecKowalskiCh1.lean:201.0-209.42
\leanok
$d_1$ is $\zeta$.
\end{theorem}
\begin{proof}
\leanok
By definition, $d_k$ is the $k$-fold Dirichlet convolution of $\zeta$. When $k = 1$, this means we are taking the convolution of $\zeta$ with itself once, which simply gives us $\zeta$. Therefore, $d_1 = \zeta^1 = \zeta$.
\end{proof}
